% !TeX root = ../discussion-review.tex
\section{Discussion}
\label{sec:discussion}

This chapter discusses the findings in relation to previous work before considering their implications and limitations.

\subsection{Discussion of the research questions}
\label{sec:discussion-rqs}

\subsubsection{Final success and recovery}

Most of the standalone questions were answered correctly on the first SQL attempt. 
Recovery contributed only a small number of final successes. 
Both models recovered on exactly five questions in each run, 
but their performance differed when execution errors occurred. 
Haiku recovered in all cases with execution errors. 
Nova failed to recover on nine and fifteen questions in the first and second runs, 
respectively (see \hyperref[tab:first-terminal-outcomes]{Table~\ref*{tab:first-terminal-outcomes}}). 
The recovery counts alone do not capture these differences.

At the same time, the models were also able to fix incorrect queries 
with the help of information returned from the database. 
However, just having the final correct answer does not tell us 
if the question was answered right the first time or took multiple attempts to get the right result. 
This is important when using a model in the real world because obtaining 
the correct result is only part of the picture. 
The recorded trajectories additionally show what issues were encountered 
and how many additional SQL attempts were needed to reach the final answer.

The absence of semantic-revision successes needs a different interpretation. 
This may be partly due to the instructions used in the experiments, 
which encouraged the agents to answer the question after a successful execution 
and did not require a separate correctness check. 
Nonetheless, the agent received the executed SQL and a preview of its results, 
so it could potentially identify a mismatch with the question. 
Thus, semantic revision was possible but was not deliberately prompted, 
and the agent received no reference-correctness feedback.

\subsubsection{What responses to errors reveal}

The models varied in their reactions to execution errors, 
and sending another query did not always fix the issue. 
The models created different SQL statements for the same questions 
and sometimes faced different errors. 
Moreover, the same error could appear in multiple attempts despite modified SQL. 
This was seen in some of Nova's cases, 
where it would change its SQL and get the same error signature on the next attempt. 
In other cases, a retry produced executable SQL, 
but the returned data still failed the reference comparison. 
Therefore, treating every SQL modification as a repair, 
or assuming that successful execution meant the problem had been resolved, 
would give a misleading account of the model's behaviour 
(\hyperref[sec:results-error-responses]{Section~\ref*{sec:results-error-responses}}).

\subsubsection{Difficulty and concentrated failures}

The gathered results show a difference between the models on hard questions.
When looking only at clean direct successes, 
Haiku still outperformed Nova (\hyperref[sec:results-difficulty]{Section~\ref*{sec:results-difficulty}}).

While the median number of attempts was one, 
there were longer strings of failures on some of Nova's queries. 
A few questions accounted for much of its SQL activity. 
These included medium-difficulty questions with long retry sequences. 
Breaking down the results by difficulty category helped to identify where performance was weaker. 
However, that was not sufficient to understand the errors behind the failures. 
The individual traces were needed to understand the errors and repeated actions behind those results.

\citet{mitsopoulou2025analysis} similarly explain that broad difficulty 
categories can obscure certain SQL challenges. 
They propose more detailed query and error analysis. 
Likewise, in this study's medium group, 
the traces showed repeated execution problems that the overall group score did not explain.

\subsubsection{Equal outcomes and repetition stability}

Final correctness does not provide a complete account of model performance. 
Comparing the models' trajectories on the same questions provided further detail. 
In the paired examples discussed in \hyperref[sec:results-paired]{Section~\ref*{sec:results-paired}}, 
both models eventually returned correct results, 
but they did not always reach them in the same way. 
In some cases, one model's first SQL attempt returned a correct result 
while the other needed to recover from an error. 
The final correctness judgment was identical, 
but the trajectory classifications revealed how the results were obtained.

Also, the comparison across runs identified another difference. 
For Nova in particular, the same question could be answered correctly 
in one run and incorrectly in the other. 
That highlights why correctness and the consistency 
of the process need to be considered separately. 
Across the two runs, Haiku showed greater consistency 
in both final correctness and trajectory classifications (\hyperref[sec:results-paired]{Section~\ref*{sec:results-paired}}).

The supplementary cases also showed 
the limits of a single measure of success. 
For example, interactive accuracy treated direct and recovered correct results equally. 
The safety cases, on the other hand, presented a different issue: 
the lack of an SQL query did not show that the model had followed the instructions. 
These situations required an examination 
of both the recorded actions and the content of the model's response 
(\hyperref[sec:results-supplementary]{Section~\ref*{sec:results-supplementary}}).
